\SourcePage{GELFAND-P000463}{unlabelled}
Рассм.
\[
K_*[x]=\operatorname{Hom}_K(K[x],K).
\]
(«дифф. операторы»).
\[
(x_*^n,x^m)=\delta_{n,m}
\]
Умножение:
\[
\underline{x_*^n\cdot x_*^m=C_{n+m}^n x_*^{n+m}}
\qquad\text{— \underline{таблица}.}
\]
Здесь $K$ не участвует.

Дост. считать $K=\mathbb Z_p$ или $\mathbb Q$.
\[
K_*[x]=\prod_{q=1}^\infty B_q,\qquad
B_q=\mathbb Z_p[y_{p^i}]/(y_{p^i}^{\,p})
\]
\Unclear{GELFAND-P000463}{product-factorization}{Слева написано $K_*[x]$. У знака произведения нижний предел $q=1$ и верхний $\infty$.}


$y_1,\ldots,y_i,\ldots,y_{p^i}$ — образующие.
\Unclear{GELFAND-P000463}{generator-powers}{В трёх местах $i$ при $p$ приподнят. Набрано $y_{p^i}$. В знаменателе у $y$ отдельно стоит верхняя степень $p$.}



\Needspace{7\baselineskip}

\SourcePage{GELFAND-P000464}{unlabelled}
\underline{30/\mbox{$\overline{\underline{\mathrm X}}$}-72г.}
\quad\underline{Сем. Гельфанда.}

\underline{И. Бернштейн}

\uuline{И.М.Г.}

\SourceFigure{GELFAND-P000464}{general-graph}{Две стрелки слева направлены к общему узлу. Из среднего узла идут стрелки влево, вправо и вниз вправо. Справа видна ломаная с наконечниками. Мелкие штрихи в узлах не прочитаны как буквы.
\par\smallskip\centering\color{black}
[DIAGRAM]}


Есть ориентированный граф. Это категория.

\sout{Р}\quad Рассм. какие л. диаграммы, где в точках — вект. пр-ва,
на стрелках отобр. (функторность не требуется).
\Unclear{GELFAND-P000464}{cancel-before-consider}{В начале зачёркнутого знака видна петля. Возможно «Р».}

\Unclear{GELFAND-P000464}{diagram-adjective}{Перед «диаграммы» предположительно «какие л.».}


Зад. — классифицировать.

\SourceFigure{GELFAND-P000464}{one-arrow}{Между двумя точками проведена слегка выгнутая стрелка вправо.
\par\smallskip\centering\color{black}
[DIAGRAM]}


Зад. о классиф. отобр.\quad Инв. лишь $\dim\operatorname{Ker}$ и $\dim\operatorname{Coker}$.
Свод. к вложению.

\SourceFigure{GELFAND-P000464}{one-loop}{Одна петля. Наконечник на её левой стороне направлен вверх.
\par\smallskip\centering\color{black}
[DIAGRAM]}


одно преобр.

\SourceFigure{GELFAND-P000464}{two-loop}{Две петли с общим левым узлом. Верхняя часть внешней дуги бледная. Возле узла видны наложенные штрихи наконечников.
\par\smallskip\centering\color{black}
[DIAGRAM]}


— зад. неразрешима (пара произв. лин. операторов). \sout{\ldots}
\Unclear{GELFAND-P000464}{cancel-before-ona}{Короткая запись перед «Она» зачёркнута и не разобрана.}


Она свод. к зад. о коммутирующих (Гельф. и Пономарёв).
\Unclear{GELFAND-P000464}{wild-name}{Первая буква наведена. Сокращение читается «Гельф.».}


\sout{Свод.\ldots}
\Unclear{GELFAND-P000464}{large-cancel-fragment}{Начало зачёркнутой строки предположительно «Свод.».}

\Incomplete{GELFAND-P000464}{large-cancel-remainder}{Остальная зачёркнутая строка не разобрана.}


\[
\left.
\begin{gathered}
\text{Пр-во и подпр-во.}\\
\text{Пр-во и 2 подпр-ва}
\end{gathered}
\right\}\quad\text{трив.}
\]
Делается:\quad пр-во и \uwave{\mbox{$3$} лин. подпр-ва}.
\Unclear{GELFAND-P000464}{three-space-word}{Начало строки предположительно «Делается».}


Инвариантов $9$, а не $8$!

\SourceFigure{GELFAND-P000464}{three-inward}{Три стрелки направлены к верхнему узлу. Их наконечники частично сливаются.
\par\smallskip\centering\color{black}
[DIAGRAM]}

(задача).

\SourceFigure{GELFAND-P000464}{four-inward}{Четыре луча сходятся к густому центральному узлу. Направления двух лучей из-за наложения штрихов не установлены.
\par\smallskip\centering\color{black}
[DIAGRAM]}


Пр-во и $4$ подпр-ва классиф. (Гельф. и Поном.) —
\sout{\ldots} уже беск. число неразл.
\Incomplete{GELFAND-P000464}{cancel-after-names}{Слово после тире зачёркнуто и не разобрано.}


\underline{Габриэль}: в каких графах конечное число неразл. систем?

Ответ: нарисовать схему Дынкина! — они и только они явл. такими графами.

Исключения: схемы с двойными (вообще кратными) стрелками.


\Needspace{7\baselineskip}

\SourcePage{GELFAND-P000465}{unlabelled}
\underline{Теорема}\quad
(\begin{tabular}[t]{@{}c@{}}Ройтов\\[-.45ex]\scriptsize из Киева\end{tabular}
или Назаровы).
\Unclear{GELFAND-P000465}{theorem-names}{Фамилии предположительно «Ройтов» и «Назаровы». Под первой строкой приписано «из Киева».}


Опис. класс схем, где есть классифик. как у Гельф. и Пономарёва.

\uwave{Связь \sout{\ldots}\ с особенностями!}\quad — см. Артин, Тюрина, Брискорн.
\Incomplete{GELFAND-P000465}{cancel-specialities}{Слово между «Связь» и «с особенностями» зачёркнуто и не разобрано.}


\begin{center}
\rule[.5ex]{18mm}{.35pt}\enspace$\cdot$\enspace\rule[.5ex]{25mm}{.35pt}
\end{center}

\uuline{И.Б.}\quad \uwave{Категория} вект. пр-в над точками и отобр. над стрелками.

(Граф фиксирован). \quad \uwave{Опред. прямая сумма}. Это $\operatorname{Cat}(\Gamma)$
($\Gamma$ — граф)

Морфизмы — отобр. над каждой точкой с коммут. диагр.

$\bullet$\qquad категория линейных пространств.

Важно: не обязательно отожд. изоморфные.
\Unclear{GELFAND-P000465}{important-line}{Сокращение перед «изоморфные» предположительно «отожд.».}


Легко видеть, что $\forall$ объект $=\sum$ неразлож. (однозн. с точн. до изом.).

$1)$\quad \uwave{Классиф. \mbox{$\Gamma$}, в к-рых конечное число неразл. объектов.}

\underline{Огранич. связными}:

\SourceFigure{GELFAND-P000465}{dynkin-a-mark}{У третьей левой вершины над горизонтальной линией виден маленький угловой штрих.
\par\smallskip\centering\color{black}
[DIAGRAM]}


\SourceFigure{GELFAND-P000465}{dynkin-d}{На отрезке между третьей и четвёртой вершинами виден небольшой наконечник влево.
\par\smallskip\centering\color{black}
[DIAGRAM]}


\SourceFigure{GELFAND-P000465}{dynkin-e}{На горизонтальной цепи шесть узлов. От третьего узла слева отходит ветвь вниз.
\par\smallskip\centering\color{black}
[DIAGRAM]}

$k=6,7,8$\qquad (слева обязательно $2$ отрезка).

Стрелки расставл. произвольно!!

$D_4$ — пр-во и $3$ подпр-ва.

Идея: нет циклов (получ. соб. значения), нет
\Unclear{GELFAND-P000465}{idea-word}{Сокращение в скобках предположительно «получ.».}


\SourceFigure{GELFAND-P000465}{forbidden-four}{Четыре отрезка сходятся к чёрной точке. На концах отдельных узлов не видно.
\par\smallskip\centering\color{black}
[DIAGRAM]}

— будет беск. число.

\SourceFigure{GELFAND-P000465}{forbidden-forks}{На обоих концах цепочки развилки. Пять узлов обозначены на самой цепочке.
\par\smallskip\centering\color{black}
[DIAGRAM]}

— тоже невозм., т.к. отожд.
\Unclear{GELFAND-P000465}{fork-explanation}{Последнее сокращение предположительно «отожд.».}


$2)$\quad \uwave{Кол-во неразл. об. не зависит от напр. стрелок.}

\uwave{Он характериз. набором размерностей}!!

надо рассм. сист. корней

$\gamma$ — полож. корень,\qquad $\alpha_i$ — пр. корни.

Каждому корню — неразл. объект
\[
\gamma=\sum n_i\alpha_i,\qquad n_i\text{ — размерности!!}
\]
\Unclear{GELFAND-P000465}{root-symbol}{Оба знака имеют нижнюю петлю и два верхних штриха. Условно набрано $\gamma$.}


Кол-во $=$ кол-ву простых корней.
\Unclear{GELFAND-P000465}{root-count}{В этой строке читается «простых корней».}


\SourceFigure{GELFAND-P000465}{a3-example}{Под цепочкой с двумя стрелками вправо стоят шесть строк размерностей.
\par\smallskip\centering\color{black}
[DIAGRAM]\par\smallskip
$\begin{array}{ccc}
1&0&0\\
0&1&0\\
0&0&1\\
1&1&0\\
0&1&1\\
1&1&1
\end{array}$}

\Unclear{GELFAND-P000465}{array-overwrite}{Средняя единица во второй строке наведена. Справа от неё виден остаток петли.}



\Needspace{7\baselineskip}

\SourcePage{GELFAND-P000466}{unlabelled}
$3)$\quad \uwave{Рассм. те \mbox{$\Gamma$}, для к-рых \mbox{$\operatorname{Cat}(\Gamma)$} разрешима.}

\SourceFigure{GELFAND-P000466}{loop-category}{Две петли имеют общий узел. В верхней петле есть короткий косой штрих, отдельно от наконечника.
\par\smallskip\centering\color{black}
[DIAGRAM]}


Неразрешим. катег. — если есть влож. $C\longrightarrow\operatorname{Cat}(\Gamma)$,
перев. неизом. объекты в неизом.

Они классиф. и соотв. расшир. схемам Дынкина!

\SourceFigure{GELFAND-P000466}{three-arrow}{Три стрелки выходят из одного узла вверх, вправо и влево вниз.
\par\smallskip\centering\color{black}
[DIAGRAM]}


\SourceFigure{GELFAND-P000466}{affine-a}{Семь точек соединены вдоль прямой. Дуга соединяет крайние точки.
\par\smallskip\centering\color{black}
[DIAGRAM]}


\SourceFigure{GELFAND-P000466}{affine-d}{Семь точек вдоль цепочки и две ветви на каждом конце.
\par\smallskip\centering\color{black}
[DIAGRAM]}


\SourceFigure{GELFAND-P000466}{affine-e}{Семь точек на горизонтали. От третьей слева отходят вниз два отрезка с точками.
\par\smallskip\centering\color{black}
[DIAGRAM]}


(появл. в теории \(\overset{\text{\sout{гр.}}}{\vphantom{x}}\) групп, порождённых отражениями.)
\Unclear{GELFAND-P000466}{cancel-in-group-comment}{Зачёркнутая приписка над строкой предположительно «гр».}


Это все графы, для к-рых зад. разрешима.

\uuline{В. Бухштабер.}

\uuline{Теор. Лазура.}
\Unclear{GELFAND-P000466}{theorem-name}{Фамилия здесь предположительно «Лазура». Возможное чтение гласной — «а».}

\[
\Lambda=\mathbb Z[a_1,\ldots,a_n,\ldots],\qquad\deg a_n=n\ge1.
\]
\underline{Док-во.}
\[
\Lambda\longrightarrow\bar\Lambda=\Lambda/(\text{разл.})
\]
\Unclear{GELFAND-P000466}{decomposables-quotient}{Над второй $\Lambda$ горизонтальная черта. В знаменателе читается «разл.».}


$\bar\Lambda$ — \uwave{универс. кольцо для}
\(\underset{\downarrow}{\boxed{\text{\scriptsize таких}}}\)
\uwave{полиномов} $h(x,y)$,
\SourceFigure{GELFAND-P000466}{such-insertion}{«таких» вписано в выноску перед «полиномов».}


\(\underset{\downarrow}{\boxed{\text{\scriptsize что}}}\)\quad
$h(x,y)=h(y,x)$,\quad $h(x,0)=0$,
\[
h(y,z)+h(x,y+z)=h(x,z)+h(y,x+z).
\]
\SourceFigure{GELFAND-P000466}{that-insertion}{«что» вписано в отдельную выноску перед условиями.}


Это когомологии:\quad $\underset{\text{(циклы)}}{dh}=0$
\qquad (в бар-конструкции для полиномов)
\Unclear{GELFAND-P000466}{cycle-caption}{Под $dh$ читается «циклы».}


\[
\operatorname{Ext}^{*,*}_{A[x]}(A,A)=H^{*,*}(A[x])
\]
\Unclear{GELFAND-P000466}{cohomology-conventions}{Под $\operatorname{Ext}$ стоит $A[x]$. В обеих частях сверху по две звёздочки.}

\[
\begin{gathered}
A[x]\longrightarrow A[x]\otimes\bar A[y]\\
x\longrightarrow x+y
\end{gathered}
\]
\Unclear{GELFAND-P000466}{normalized-coproduct}{Под первой стрелкой написано $x\longrightarrow x+y$.}


Когом. счит. с пом. резольвент.

Бар-констр. — одна из резольвент.

\[
\mathcal A:\quad
A[x]\xrightarrow{\Delta}A[x]\otimes\bar A[y]
\longrightarrow A[x]\otimes\bar A[y]\otimes\bar A[z]\longrightarrow\ldots
\]
$\bar A[y]$ — многочлены без своб. члена.

Получ. ацикличный комплекс, к-рый прим. для вычисл. когомологий
\[
(\text{ставим }\underset{h(x,y)}{\text{всюду}}\text{ — }
\tilde a\Rightarrow a\otimes1+1\otimes a)=\Delta a).
\]
\Unclear{GELFAND-P000466}{primitive-condition}{Над $a$ волнистая черта. Следующий знак предположительно $\Rightarrow$. Видны обе закрывающие скобки.}


Надо классиф. циклы в конкретной резольвенте.
\[
0\longrightarrow B\longrightarrow Z\longrightarrow H^2(A[x])\longrightarrow0
\]
\Unclear{GELFAND-P000466}{exact-sequence}{После $B$ написано $Z$, у $H$ верхний индекс $2$.}


$x^n\longrightarrow(x+y)^n-x^n-y^n$\qquad (это границы)


\Needspace{7\baselineskip}

\SourcePage{GELFAND-P000472}{unlabelled}
\SourceFigure{GELFAND-P000472}{top-product-diagram}{Узел $C_1$ густо наведён. От $C_i$ к $A$ идёт прерывистая стрелка, у стрелки от $C_{i+j}$ к $C_j$ один разрыв.\par\smallskip\centering
[DIAGRAM]}


Получили гомоморфизм $C_{i+j}\longrightarrow A_0$.
\Unclear{GELFAND-P000472}{homomorphism-verb}{Первое слово читается «Получили», предпоследняя буква сближена с последней.}


Получится
\[
A\underset{\text{сопост.}}{\rightsquigarrow}H(A)
\qquad\text{(кольцо)}
\]
$H(A)$ вообще говоря некоммутативно (!) по настоящему.
На этом основана классифик. многозн. групп.
\Unclear{GELFAND-P000472}{multivalued-groups}{После «классифик.» читается «многозн.».}

(В топологии до сих пор встр. только коммут. кольца — там всегда алгебры Хопфа).

\uline{Станд. резольвента или \mbox{$B$}-резольвента.}
\[
\overline A=\sum_{n>0}A_n
\]
\[
C_1=A\underset{A_0}{\otimes}\overline A
\qquad\text{В дальнейшем }A_0\text{ опускается}
\]
(Исх. кольца коммутативны!)
\[
\begin{gathered}
C_s=A\otimes(\overline A)^s.\\
d_{s+1}\colon A\otimes(\overline A)^{s+1}
\longrightarrow A\otimes(\overline A)^s
\end{gathered}
\]
\[
\begin{aligned}
d_s\bigl(a_0[a_1|a_2|\ldots|a_s]\bigr)
={}&(-1)^{\varepsilon(0)}a_0a_1[a_2|\ldots a_s]\\
 &+\sum(-1)^{\varepsilon(n)}a_0[a_1|\ldots|a_na_{n+1}|\ldots a_s]
\end{aligned}
\]
\Unclear{GELFAND-P000472}{bar-differential-indices}{Два индекса в произведении наведены. Читается $a_na_{n+1}$.}


$\varepsilon(n)$ — некоторые ф-ии от номера,
\sout{Для чётной градуир. \mbox{$\varepsilon(n)=n$}.}
\par\uline{\mbox{$\varepsilon(n)=n$}, если все элем. чётные.}
\SourceFigure{GELFAND-P000472}{even-sign-cancellation}{Условие о чётных элементах вписано ниже зачёркнутой строки и отделено чертой.}


\[
\operatorname{Hom}_A(C_s,A_0)\simeq(\overline A^{\,*})^s,
\qquad \overline A^{\,*}=\operatorname{Hom}_{A_0}(\overline A,A_0)
\]
\[
d_s^*\colon[\alpha_1|\alpha_2|\ldots|\alpha_{s-1}]
=\sum(-1)^{\varepsilon(n)}
[\alpha_1|\ldots|\overline\Delta\alpha_n|\ldots|\alpha_{s-1}].
\]
\[
\varepsilon(n)=n,\quad\text{если }\deg\alpha_i=2k_i
\]

\uline{Опред. \mbox{$\overline\Delta$}.}
\SourceFigure{GELFAND-P000472}{lower-coproduct-box}{У верхней стрелки явное правое остриё; левый конец наведён и также похож на остриё.\par\smallskip\centering
[DIAGRAM]}


На самом деле
\[
\operatorname{Im}\overline\Delta\subset\overline{A^*}.
\]
\Unclear{GELFAND-P000472}{lower-reduced-image}{Черта справа проходит над $A$ и звёздочкой.}


Отсюда ещё одно определение умножения,
\[
\begin{aligned}
C_s^*&\qquad f=[\alpha_1|\ldots|\alpha_s]\\
C_{s'}^*&\qquad g=[\alpha'_1|\ldots|\alpha'_{s'}]
\end{aligned}
\qquad
f\cdot g=[\alpha_1|\ldots|\alpha_s|\alpha'_1|\ldots|\alpha'_{s'}]
\]
\Unclear{GELFAND-P000472}{bottom-prime-index}{Последний индекс у $g$ читается $s^{\prime}$. В записи произведения сильно наведена граница после первого штрихованного члена.}



\Needspace{7\baselineskip}

\SourcePage{GELFAND-P000473}{unlabelled}
\uline{20/\mbox{$\overline{\underline{\mathrm{XI}}}$}-72.}
\qquad\uline{Сем. Гельфанда.}

\uuline{Костюченко.}\qquad\uuline{Один пример.}
\SourceFigure{GELFAND-P000473}{date-heading}{Месяц обозначен римской цифрой. Имя и название подчёркнуты дважды.}

\[
\begin{gathered}
v''+2a\lambda v'+b\lambda^2v=0\\
v(0)=v(\pi)=0.
\end{gathered}
\qquad\text{в }L_2[0,\pi].
\]

\SourceFigure{GELFAND-P000473}{pde-bracket}{Под уравнением линия с короткими загнутыми вверх концами.\par\smallskip\centering
[DIAGRAM]}

\[
\begin{gathered}
0<x<\pi\\
0\leqslant t<+\infty
\end{gathered}
\]
\Unclear{GELFAND-P000473}{top-time-inequality}{Второй знак неравенства густо наведён. Условное чтение $0\leqslant t<+\infty$.}

\uline{Пусть эллипт. ур-е.}
\Unclear{GELFAND-P000473}{elliptic-heading}{Первое слово читается «Пусть».}


\[
\underline{a^2<b};\qquad\underline{b>0}.
\]
\[
v_k(x)=e^{\gamma kx}\sin kx,\qquad
\gamma=\frac{-a}{\sqrt{b-a^2}};\qquad k=\pm1,\ \pm2,\ \ldots .
\]
\SourceFigure{GELFAND-P000473}{gamma-bracket}{Под формулой для $\gamma$ рамочная черта с загнутыми концами.}


1. Система полна, но неизв., как находить коэффициенты.
Это вопрос о нахожд. биортог. системы.
\Unclear{GELFAND-P000473}{upper-system-correction}{Окончание слова «системы» наведено. После точки отдельный зачёркнутый знак.}

(нах. в терминах обр. ф-и к интегр. Кристоффеля–Шварца)

(\uline{биортог. система есть})

2. \uline{Это не базис!} (не любой вектор разлаг. в сход. ряд по этой сист.)

\sout{Даже}\quad\uline{Более того:} если есть разл. в $L^2$,
то предел — аналит. ф-я! — и вып. беск. число краев. условий (верно и обратное).
\Unclear{GELFAND-P000473}{middle-cancelled-even}{Зачёркнутое слово предположительно «Даже».}


Нельзя суммировать методом Абеля 1-го порядка.

\uline{если}
\[
\sum_{\lambda_k>0}e_k e^{\lambda_k t}v_k(x)
\]
\Unclear{GELFAND-P000473}{middle-abel-coefficient}{Коэффициент густо наведён; условное чтение $e_k$.}

\Incomplete{GELFAND-P000473}{middle-abel-cancelled}{Два коротких фрагмента после $v_k(x)$ густо зачёркнуты. Чтение не установлено.}

им. предел при $t\to0$,
то \uwave{ф-я предельная} аналитична.
\SourceFigure{GELFAND-P000473}{limiting-function-line}{Волнистая черта проходит над «ф-я» и под «предельная».}


(Дискредитация метода Фурье для эллипт. ур-я)

При нек-рых неразд. условиях можно решать.
\Unclear{GELFAND-P000473}{boundary-abbreviation}{Перед «условиях» читается сокращение «неразд.».}


\Incomplete{GELFAND-P000473}{lower-cancelled-marks}{Перед «Причина» видны две зачёркнутые пометы. Чтение не установлено.}

\uline{Причина плохого повед.:} резольвента $\exp$ растёт в углах
\SourceFigure{GELFAND-P000473}{resolvent-angle-diagram}{Две пересекающиеся линии, прерывистые штрихи слева, точка на пересечении, кружок и стрелка справа.\par\smallskip\centering
[DIAGRAM]}


На $[0,\pi-\varepsilon]$ переопределена.

На $[0,\pi+\varepsilon]$, $0<\varepsilon<\pi$ она остаётся полной !!! (И. Бернштейн)
\Unclear{GELFAND-P000473}{epsilon-lower-bound}{Нижняя граница густо наведена и условно набрана как ноль.}


На $[0,2\pi]$ возникает много ортог. ф-ий.

\uline{Задача теории целых ф-ий.}
\[
\int_0^\pi f(x)e^{\gamma kx}\sin kx\,dx=0
\quad\Longrightarrow\quad
F(\lambda)=\int_0^\pi f(x)e^{\gamma\lambda x}\sin\lambda x\,dx
\]
$F(\lambda)$ целая ф-я $\exp$ типа. Она 0 в целых точках
$F(n)=0$, $n>0$.

\uline{Надо учесть специфику.}

На $[0,\pi-\varepsilon]$ срабатывает тип и индик. диаграмма.

Единственность — ?

Неизвестно, что будет, если заменить $k$ на $\lambda_k$.


\Needspace{7\baselineskip}

\SourcePage{GELFAND-P000474}{unlabelled}
\uuline{И. М. Гельфанд}\qquad\uuline{о раб. Арнольда.}
\SourceFigure{GELFAND-P000474}{gelfand-arnold-heading}{Имя и тема дважды подчёркнуты.}

\[
(\mathbb R^n,0)\longrightarrow(\mathbb R,0)
\]
$f(x)=0$ — гиперпов.

Как ведёт себя гиперпов. в окр. $0$ — ?

$f\in C^\infty$;\quad $f(0)=0$, $df(0)=0$. (крит. точка).

(На самом деле рассм. $\mathbb C^n\longrightarrow\mathbb C$)

Эквив. гиперпов. — отлич. заменой переменных
\[
(\mathbb R^n,0)\longrightarrow(\mathbb R^n,0)\qquad
\text{(диффеом.)}.
\]

\uline{Простая гиперпов.:}\qquad\uline{\mbox{$f(x)=0$}}

если пошевелить, то будет конечное число орбит
в струях пор. $N$ (оценка числа орбит не зависит от номера $N$).

Перечисление простых орбит:
\SourceFigure{GELFAND-P000474}{simple-chain-graphs}{Слева две цепочки. В конце нижней цепочки развилка.\par\smallskip\centering
[DIAGRAM]}

\SourceFigure{GELFAND-P000474}{exceptional-graphs}{Справа цепочка с нижним ответвлением у третьей вершины. Последние звенья прерывистые.\par\smallskip\centering
[DIAGRAM]}


Далее: Большинство орбит простые (как правило, орбита-простая).
\Incomplete{GELFAND-P000474}{middle-parenthetic-cancelled}{После открывающей скобки зачёркнут короткий фрагмент. Чтение не установлено.}

\uline{Непростые коразмерности 6.}

\[
x_1^{n+1}+x_2^2+\cdots\,\ast=0
\]
— соотв. 1-му типу.
\Unclear{GELFAND-P000474}{middle-simple-equation}{Первая буква в показателе $n+1$ наведена. После многоточия виден отдельный знак, похожий на звёздочку. Под формулой рамочная черта.}


\uline{Версальная деформация.}
\[
\begin{gathered}
F(x,\lambda)=0,\qquad F(x,0)=f(x),\\
\lambda=(\lambda_1,\lambda_2,\ldots,\lambda_k)
\end{gathered}
\]
Всякая другая деф. эквив. деф., индуц. из этой след. образом:
\[
G(x,\mu)=F(X(x,\mu),\Psi(\mu))
\quad\text{и}\quad X(y,0)=y,\quad\Psi(0)=0
\]
\Unclear{GELFAND-P000474}{middle-induced-variables}{Первые аргументы $G$ и $X$ густо наведены и условно читаются как $x$. Справа ясно написано $X(y,0)=y$.}


\uline{Правило построения верс. сем-ва:}
Берём орбиту и мн-ие, трансверс. к орбите.
\SourceFigure{GELFAND-P000474}{transverse-sketch}{Изогнутая линия и короткий поперечный отрезок с точкой пересечения.\par\smallskip\centering
[DIAGRAM]}

{\centering\rule{1.8em}{.4pt}\quad\textperiodcentered\quad\rule{2.4em}{.4pt}\par}

\[
F(x,\lambda)=f(x)+
\sum a_i\left.\frac{\partial F}{\partial\lambda_i}\right|_{\lambda_i=0}
=f(x)+\sum c_i\varphi_i(x)+O(\lambda^2).
\]
\Unclear{GELFAND-P000474}{lower-taylor-coefficient}{Множитель первого слагаемого суммы условно читается $a_i$. Под чертой подстановки написано $\lambda_i=0$.}


\uuline{\mbox{$\varphi_i$} зад}\quad\uline{Инфинитезимальную версальную деформацию:}
если $\forall$ ф-я $a(x)$ предст. в виде
\Unclear{GELFAND-P000474}{lower-definition-heading}{Первый фрагмент густо наведён и дважды подчёркнут. Условное чтение $\varphi_i$ «зад».}

\[
a(x)=\sum\underline{\underline{\frac{\partial f}{\partial x_i}}}\cdot h_i(x)
+\sum_{i=1}^{k}c_i\varphi_i(x)
\]
$a(x)$ задаёт инфинитезим. верс. деформ.

($\varphi_i$ порожд. лок. кольцо глад отобр.)

\Incomplete{GELFAND-P000474}{bottom-cancelled-taylor}{Перед скобками с записью о Тейлоре видна отдельная зачёркнутая запись. Чтение не установлено.}

(из разл. по Тейлору той ф-ии)

Важно: факторкольцо по \uline{якобиеву идеалу}
$\left(\left\{\dfrac{\partial f}{\partial x_i}\right\}\right)$.

\uline{Фактор} — конечномерное пр-во размерности $k$.

{\centering\rule{2.9em}{.4pt}\quad\textperiodcentered\quad\rule{3.4em}{.4pt}\par}


\Needspace{7\baselineskip}

\SourcePage{GELFAND-P000475}{unlabelled}
\[
\underline{\underline{f=z_1^{k+1}+z_2^2+\cdots+z_n^2.}}
\]

Тогда
\[
\frac{\partial f}{\partial z_1}=z_1^k,\qquad
\frac{\partial f}{\partial z_2}=z_2,\ \ldots,
\]
\Unclear{GELFAND-P000475}{top-derivatives}{Справа написаны одночлены без числовых множителей.}

отсюда факторкольцо — пр-во размерности $k$.

\[
\varphi_i=1,z_i,z_i^2,\ldots,z_1^{k-1}
\]
\Unclear{GELFAND-P000475}{top-phi-indices}{У первых двух степеней после $1$ видны индексы, похожие на $i$; у последней — $1$.}


Верс. деформация:
\[
\underline{\lambda_0+\lambda_1z_1+\cdots+\lambda_{k-1}z_1^{k-1}
+z_1^{k+1}+z^2}
\quad\text{— \uline{верс. деформация}}
\]
\SourceFigure{GELFAND-P000475}{versal-bracket}{Под многочленом рамочная черта. Последнее слагаемое написано $z^2$.}


{\centering\rule{2.8em}{.4pt}\quad\textperiodcentered\quad\rule{3.2em}{.4pt}\par}
\uline{\mbox{$F(z,\lambda)=0$}.}\qquad
Бифуркац. диаграмма — мн-во тех $\lambda$, где остаются особ. точки (в верс. деформации).

\SourceFigure{GELFAND-P000475}{theorem-star}{Слева от теоремы четыре отрезка пересекаются в одной точке.\par\smallskip\centering
[DIAGRAM]}

\uuline{Теор.:} на откр.
\(\underset{\text{\scriptsize\tikz[baseline=(label.base)]\node[draw,ellipse,inner sep=1pt] (label) {бифурк. диагр.};}}{\text{мн-ве}}\)
уже будут простые особ. точки (морсовские).
\SourceFigure{GELFAND-P000475}{theorem-bifurcation-insertion}{Сокращение «бифурк. диагр.» вписано в выноску под «мн-ве».}


\[
z_1^{k+1}+z_2^2=0
\qquad\Vert\qquad
\begin{gathered}
\text{бифурк. диагр. этой точки —}\\[-.2ex]
\longrightarrow\text{ Ласточкин хвост}.
\end{gathered}
\]
\[
\begin{gathered}
z_1^{k+1}+z_2^2+\lambda_{k-1}z_1^{k-1}+\cdots+\lambda_0\\
z_2^2=-p_{k+1}(z_1);\qquad z_2=\sqrt{-p_{k+1}(z_1)}.
\end{gathered}
\]
особ. точки\quad $(z_2,P'(z_1))=0$, т. е.
\[
\underline{P(z_1)=0\ \text{и}\ P'(z_1)=0}
\]
($P$ имеет кратные корни)
\SourceFigure{GELFAND-P000475}{multiple-root-bracket}{Оба равенства обведены снизу общей рамочной чертой.}


Рассм. доп. к бифурк. диаграмме и его фунд. группу.
Она есть группа кос группы Коксетера.

Рассм. ур-е
\[
\underline{P(z,\lambda)=0,\quad |z|<\varepsilon}
\qquad\text{\uline{\mbox{$\lambda$} такое, что нет особенностей}}
\]
\Unclear{GELFAND-P000475}{middle-polynomial-name}{Первая буква густо наведена. Условное чтение $P$.}


гомологии этого пр-ва $\ne0$ лишь в \sout{разм.}
\Incomplete{GELFAND-P000475}{lower-homology-cancelled}{После зачёркнутого «разм.» видны остатки записи. Полное чтение не установлено.}

половинной разм. $n-1$.

Оно есть букет из $\mu$ сфер.
\SourceFigure{GELFAND-P000475}{sphere-bouquet}{Четыре петли с общей точкой изображают букет.\par\smallskip\centering
[DIAGRAM]}


\Incomplete{GELFAND-P000475}{lower-left-cancelled}{Запись слева от $H^{n-1}$ густо зачёркнута. Чтение не установлено.}

$H^{n-1}$ имеет образующие $e_1\ e_2\ldots\ e_\mu$.

Индекс пересеч. опред. скалярное произвед. $(e_i,e_k)$.

\[
n-1\equiv2\pmod4\qquad\Vert\qquad
\begin{gathered}
\text{отриц. опред.}\\
\text{полуопред.}\\
\text{неопред.}
\end{gathered}
\qquad\underline{e_i^2=-2}\quad\text{(всегда)}
\]
\Incomplete{GELFAND-P000475}{lower-signature-cancelled}{Между «полуопред.» и «неопред.» одна зачёркнутая строка. Чтение не установлено.}


\uwave{монодромия}: предст. фунд. гр. в
\Incomplete{GELFAND-P000475}{lower-monodromy-cancelled}{После «в» зачёркнута запись; дальше многоточие. Чтение записи не установлено.}

\ldots

{\centering\rule{4em}{1pt}\quad\textbullet\quad\rule{4em}{1pt}\par}

\sout{Атья}
\Unclear{GELFAND-P000475}{bottom-cancelled-name}{Зачёркнутое имя предположительно «Атья».}

\uline{Есть раб. Атья, Ботта,\ldots:} как избавиться от производных в выкладках Синг.?
Как, сидя в одной точке, изучить топологию мира? Нахожд. явного вида асимптотики для ур-й 2-го порядка (сообр. И. М. Гельфанда, В. И. Арнольда и др.)

(будет в следующем семестре)

А. Шнирельман о работах Маслова.


\Needspace{7\baselineskip}

\SourcePage{GELFAND-P000476}{unlabelled}
\uuline{В. М. Бухштабер.} (продолжение).

\[
\begin{gathered}
A\overset{\text{аддит.}}{=}\Lambda[[X]],\qquad
X_n\text{ — аддит. образующие}\\
X_n\cdot X_m=C_{n+m}^n X_{n+m}
\end{gathered}
\]
\Unclear{GELFAND-P000476}{additive-equality}{Под «аддит.» две короткие параллельные черты, читается знак равенства.}


\[
\begin{gathered}
\operatorname{Ext}_A^{2,*}(\Lambda,\Lambda)\\[-.4ex]
\rotatebox[origin=c]{90}{$=$}\\[-.4ex]
\operatorname{Ker}d_3^*/\operatorname{Im}d_2^*
\end{gathered}
\qquad\text{в терминах }B\text{-конструкции}
\]

$1^\circ$.
\[
d_3^*\colon\overline A^*\otimes\overline A^*
\longrightarrow\overline A^*\otimes\overline A^*\otimes\overline A^*
\]

\[
A^*=\operatorname{Hom}_{\Lambda}^{c}(\Lambda[[X]],\Lambda)
\overset{\text{адд.}}{\simeq}\Lambda[x]
\]
\Unclear{GELFAND-P000476}{top-dual-hom}{У $\operatorname{Hom}$ верхний индекс читается $c$, нижний — $\Lambda$.}

\uuline{компактное} (ненулев. знач. на конечн. числе).
\SourceFigure{GELFAND-P000476}{compact-annotation}{Пояснение «компактное» написано под $\operatorname{Hom}$ и соединено с ним штрихом.}


\[
\begin{gathered}
\Delta\colon A^*\longrightarrow A^*\otimes A^*\\
\Delta x_n\longrightarrow x_n\otimes1+1\otimes x_n+\sum d_{i,j}x_i\otimes x_j^*
\end{gathered}
\]
\Unclear{GELFAND-P000476}{top-delta-star}{После $x_j$ видна звёздочка. В следующей формуле для $\Delta x_n$ звёздочки нет.}


\[
\begin{aligned}
(\Delta x_n,x_i\otimes x_j)&=d_{i,j}\\[-.4ex]
\qquad\rotatebox[origin=c]{90}{$=$}&\\[-.4ex]
(x_n,x_i\cdot x_j)&=(x_n,C_{i+j}^{i}x_{i+j})=C_{i+j}^{i}\delta_{n,i+j}
\end{aligned}
\]
\SourceFigure{GELFAND-P000476}{pairing-layout}{Две строки скалярного спаривания соединены вертикальным знаком равенства.}


\[
\Delta x_n=x_n\otimes1+1\otimes x_n+\sum C_n^i x_i\otimes x_{n-i}
\]
И положим\quad $x_n=x^n$
\[
\begin{gathered}
\Delta x^n=(x\otimes1+1\otimes x)^n\\
d_3^*\colon \overline{\Lambda[x,y]}\longrightarrow\overline{\Lambda[x,y,z]}
\end{gathered}
\]

\[
d_3\Biggl(\underset{
  \substack{\rotatebox[origin=c]{90}{$=$}\\ \sum h_{i,j}x^iy^j}
}{h(x,y)}\Biggr)=\sum h_{ij}(\overline\Delta x^i)z^j
-\sum h_{i,j}x^i\overline\Delta y^j=
\]
\Unclear{GELFAND-P000476}{middle-differential-star}{В вычислении перед $h(x,y)$ стоит $d_3$ без звёздочки. В предыдущей строке написано $d_3^*$.}


\[
\overline\Delta x^n=
\underset{\substack{
 \rotatebox[origin=c]{90}{$=$}\\
 (\Delta(x))^n\\
 \rotatebox[origin=c]{90}{$=$}\\
 (x+y)^n
}}{\Delta x^n}-x^n-y^n.
\]
\SourceFigure{GELFAND-P000476}{coproduct-identifications}{Под $\Delta x^n$ вертикальными знаками равенства соединены $(\Delta(x))^n$ и $(x+y)^n$.}


\[
=h(x+y,z)-\mbox{\sout{\ensuremath{h(x,z)}}}-h(y,z)
-h(x,y+z)+h(x,y)+\mbox{\sout{\ensuremath{h(x,z)}}}.
\]
\SourceFigure{GELFAND-P000476}{cancelled-cocycle-terms}{Оба вхождения $h(x,z)$ перечёркнуты. Знаки минус и плюс перед ними видны.}


\uuline{\mbox{$\operatorname{Ker}d_3^*$}:}
\[
\boxed{h(x+y,z)+h(x,y)=h(x,y+z)+h(y,z)}
\]
Эти циклы обозн. $L$.
\SourceFigure{GELFAND-P000476}{kernel-relation-box}{Равенство заключено в прямоугольную рамку.}


Это и есть лазаровское соотношение.
\Unclear{GELFAND-P000476}{lazard-name}{Перед «соотношение» читается «лазаровское».}


\[
\underline{\operatorname{Im}d_2^*}=\Lambda\{d_2^*x^n\}
\]
\[
0\longrightarrow\operatorname{Im}d_2^*\longrightarrow L
\longrightarrow\operatorname{Ext}_A^{2,n}(\Lambda,\Lambda)\longrightarrow0
\]

\uline{Счёт группы \mbox{$\operatorname{Ext}$}.}

В ур-ии не фигурирует $\Lambda$. Поэтому дост. счит. 2 случая:

1)\quad $\Lambda$ — $Z_p$-модуль ($p$ — простое)

2)\quad $\Lambda=Q$


\Needspace{7\baselineskip}

\SourcePage{GELFAND-P000490}{unlabelled}
\uline{Теор. Милнора — Кервера.}\quad Все $M_1$ получ. друг из друга перестр. Морса.

(необх. и дост.)

Сущ. ли там мн-зие, гомот. эквив. $M$.

$P_n$ — группа препятствий
\Unclear{GELFAND-P000490}{obstruction-cancellation}{После «препятствий» зачёркнуто слово. Чтение не установлено.}

к перестройкам до гомотопич. эквив.
\[
\underline{\sigma_0:M_1\longrightarrow M}
\]
\uline{Послед. точна.}
\[
?\longrightarrow hS(M^n).
\]
\SourceFigure{GELFAND-P000490}{exact-question}{На месте первого члена написан вопросительный знак.}


\[
\begin{aligned}
\alpha_1.&\quad\{N_1\overset{f_1}{\longrightarrow}M\},\\
\alpha_2:&\quad\{N_2\overset{f_2}{\longrightarrow}M\}.
\end{aligned}
\]
\SourceFigure{GELFAND-P000490}{two-structures}{После $\alpha_1$ стоит точка, после $\alpha_2$ — двоеточие.}


$\xi_1\in\phi(M)$ постр. по $\alpha_1$.
\SourceFigure{GELFAND-P000490}{group-sign}{Знак перед $(M)$ передан как $\phi$.}


\SourceFigure{GELFAND-P000490}{thom-pair}{Две стрелки к комплексам Тома. Между правыми членами два вертикальных штриха.\par\smallskip
{\centering
[DIAGRAM]\par}}

\Unclear{GELFAND-P000490}{bundle-symbol}{Волнистая черта может быть частью $\tau_M$ либо отдельной тильдой над $\tau$. Условно набрано $\tau_M$.}


Получ. задача:
\Unclear{GELFAND-P000490}{novikov-cancellation}{Между «задача» и «дисс.» зачёркнута запись. Чтение не установлено.}

дисс. Новикова:
\SourceFigure{GELFAND-P000490}{homotopy-target}{Под правым концом стрелки стоят $\cup$ и $M$.\par\smallskip
{\centering
[DIAGRAM]\par}}


Пленка:
\[
\begin{aligned}
F\big|_{S^L\times0}&=\beta_1\\
F\big|_{S^L\times1}&=\beta_2
\end{aligned}
\]
\SourceFigure{GELFAND-P000490}{endpoint-restrictions}{Два ограничения $F$ при $0$ и $1$.}

\Unclear{GELFAND-P000490}{last-line-cancellation}{Перед квантором затёрт знак. Чтение не установлено.}


$\exists$ ли среди пленок $h$-кобордизмы?


\Needspace{7\baselineskip}

\SourcePage{GELFAND-P000491}{unlabelled}
\uline{19/\mbox{$\overline{\underline{\mathrm{III}}}$}-73г.}
\qquad\uline{Сем Гельфанда}
\Unclear{GELFAND-P000491}{date-cancellation}{Перед датой зачёркнута короткая запись. Возможно, цифры.}


\uuline{Арнольд.}\quad Рассм. росток ф-и $f:(\mathbb C^n,0)\longrightarrow(\mathbb C,0)$.
$0$ — крит. точка.
\SourceFigure{GELFAND-P000491}{annotation-icon}{Над строкой виден жёлтый значок выноски с тремя чёрными штрихами. Текст выноски в кадре не раскрыт.}


Рассм. пр-во $k$-струй. Действ. гр. $k$-струй дифф.

$j^k f$\qquad $k$-достат. (след. не влияют)

\uline{\mbox{$m$} — число модулей.}
\uuline{Теор..}\quad $m=0\Longrightarrow$ список
\[
\begin{array}{ll}
A_k & x^{k+1}\\
D_k & x^2y+y^{k-1}\\
E_6 & x^3+y^4\\
E_7 & x^3+xy^3\\
E_8 & x^3+y^5
\end{array}
\]
\SourceFigure{GELFAND-P000491}{ade-table}{В списке усилены некоторые штрихи, в том числе у $E_6$ и $E_8$.}


Стаб. экв. \quad $\oplus$ с невыр. квадр. формами.

Квадр. форма полож. определена.


\Needspace{7\baselineskip}

\SourcePage{GELFAND-P000492}{unlabelled}
Аналит. ф-я с точн. до анал. зам. перем. однозн. опр. зам. перем.

$0$-модульная особ. — после изм. коэфф. $\exists$ лишь конечное число вар.
\Unclear{GELFAND-P000492}{modality-ending}{Окончание слова исправлено. В строке условно читается «модульная». Возможно «модульных».}


{\centering\rule{1.4em}{.4pt}\quad\textperiodcentered\quad\rule{1.6em}{.4pt}\par}

$f=\varepsilon$\quad лин. уровня:
\uline{Пересеч. с шаром} — мн-зие с краем чётной разм.
\[
f^{-1}(\varepsilon)\cap B_\delta
\]
\SourceFigure{GELFAND-P000492}{fibre-drawing}{Две замкнутые кривые соединены вогнутыми боковыми линиями.\par\smallskip
{\centering
[DIAGRAM]\par}}


Индекс пересеч.
\Unclear{GELFAND-P000492}{intersection-cancellation}{После «пересеч.» зачёркнута запись. Чтение не установлено.}

вещ. гомол. средней разм. (целочисл.)

Она
$\overset{\text{\scriptsize (отр. опр.)}}{\text{знакоопр.}}$
в точности в перечисл. случаях,
если только $n\equiv3\ (\mathrm{mod}\ 4)$.
(иначе доб. лишние перем.) — теор. Тюрина.
\SourceFigure{GELFAND-P000492}{negative-insert}{«(отр. опр.)» вписано в выноску после «знакоопр.».}

\[
f^{-1}(\varepsilon)\cap B_\delta\approx
\underbrace{S^{n-1}\vee\ldots\vee S^{n-1}}_{\mu}.
\]
$\mu=\dim$ лок. кольца.

\SourceFigure{GELFAND-P000492}{dynkin-diagram}{Схема $E_6$. Возле крайней левой вершины написано $-2$, под первым ребром — $1$.\par\smallskip
{\centering
[DIAGRAM]\par}}

\Unclear{GELFAND-P000492}{graph-cancellation}{Справа от $E_6$ зачёркнута запись. Чтение не установлено.}


Кв. форма зад. сх. Дынкина.

\fbox{$m=1$.} — тоже находится список (длинный).

Первая:
[DIAGRAM]

\begin{center}
[DIAGRAM]
\end{center}

[DIAGRAM]
$\Longrightarrow$ всегда $m=1$.
\Unclear{GELFAND-P000492}{modality-inequality}{Знак перед $1$ усилен штрихами. Условно набрано $\leqslant$. В заключении написано $m=1$.}


У них у всех
\[
\mu_+\leqslant1\Longrightarrow\mu_0=1\qquad(\text{всегда вырожд.})
\]
\uuline{\mbox{$\mu_+=1$}.}\quad
$axyz+x^k+y^\ell+z^m$ — \uline{3-х инд. серии} $T_{k\ell m}$.

\uuline{А. С. Мищенко.}

\uline{Сулливан:} локализация, периодичн. и \ldots\ Галуа-
\[
K(X)/\text{\mbox{$J$}-эквив.}=J(X)\quad\text{— \uline{конечн. группа}}.
\]
Адамс: она опис. в терм. когом. опер. $K$-теории $\psi^k$.
\[
\psi^k(\xi)=\xi^k\quad\text{и \uwave{аддитивны}}
\quad\Longrightarrow\quad\text{этим они однозн. опред.}
\]
$(\dim\xi=1)$\qquad $k=2,3,\ldots$

Итак $J(X)$ опр. группой $K(X)$ как модулем над этими опер..

$J$-эквив. — послойная гомот. экв..
\SourceFigure{GELFAND-P000492}{fibrewise-diagram}{Стрелка $f$ соединяет $\xi$ и $\eta$. От обеих букв проведены стрелки к $X$.\par\smallskip
{\centering
[DIAGRAM]\par}}

$f$ — собств. \quad $\deg f=1$ в слоях.

Было док. (Адамс, Атья):
\[
\xi\overset{J}{\sim}\text{трив.}.
\]

$X$ — конечн. комплекс.
\raisebox{.45ex}{\scriptsize (тогда $\exists N\in\mathbb Z$:)}
\SourceFigure{GELFAND-P000492}{integer-quantifier}{После «комплекс» вписано «тогда $\exists N\in\mathbb Z$:».}


$\forall$ $J$-трив. рассл. $\xi$
\raisebox{.45ex}{\scriptsize ($\forall k$)}
$\exists$
\SourceFigure{GELFAND-P000492}{k-quantifier}{Между $\xi$ и $\exists$ вписано $\forall k$.}

\Unclear{GELFAND-P000492}{bundle-cancellation}{Перед «рассл.» зачёркнуто слово. Чтение не установлено.}

рассл. $y$:
\Unclear{GELFAND-P000492}{formula-cancellation}{После двоеточия зачёркнута запись. Чтение не установлено.}

\begin{center}
[DIAGRAM]
\end{center}
\SourceFigure{GELFAND-P000492}{adams-equality-bracket}{Равенство обведено снизу и с боков.}


\uline{Обратное — гипотеза Адамса:}
$\forall$ кон. компл. $X$ $\exists N$; $\forall k$
\[
k^N(\psi^k-1)\xi\overset{J}{\sim}0.
\]


\Needspace{7\baselineskip}

\SourcePage{GELFAND-P000493}{unlabelled}
\uline{Идея.}
\SourceFigure{GELFAND-P000493}{universal-bundle-diagram}{Слева написано $\xi=f^*(\xi)$, справа $\xi_n$.\par\smallskip
{\centering
[DIAGRAM]\par}}

\Unclear{GELFAND-P000493}{pullback-star}{После $f$ густо наведён знак, условно набранный как звёздочка.}


\SourceFigure{GELFAND-P000493}{pullback-square}{Вторая схема с $f^*(\eta)$ и $\eta$.\par\smallskip
{\centering
[DIAGRAM]\par}}

\[
\psi^k(f^*(\eta))=f^*\psi^k(\eta).
\]
Отсюда опер. $\psi^k$ опр. заданием рассл. на \uuline{\mbox{$BU(n)$}}.
$(\psi^k(\xi_n))$.

По $\psi^k$ строится отобр.
$BU(n)\longrightarrow BU(\infty)$ (неизв. размерность).

$\psi^k(1)=1\Longrightarrow$ получ., что $\exists$ послед. отобр.
\SourceFigure{GELFAND-P000493}{stabilization-diagram}{Над $BU(n)$ написано $\xi_n+1$.\par\smallskip
{\centering
[DIAGRAM]\par}}


Итак, $\psi^k$ зад. одним отобр.
\begin{center}
[DIAGRAM]
\end{center}
\SourceFigure{GELFAND-P000493}{single-arrow}{Стрелка обведена снизу двойной чертой с загнутыми концами.}


\uwave{Пусть есть когом. опер.}
$BU(\infty)\overset{A}{\longrightarrow}BU(\infty)$,
и пусть она сохр. разм.
\SourceFigure{GELFAND-P000493}{rank-square}{Квадрат с $A$ на обеих горизонтальных стрелках.\par\smallskip
{\centering
[DIAGRAM]\par}}


\uwave{сохр. размерность.}
\uwave{Их не бывает, но если бы они были, то можно было бы док., что}
$A\xi\overset{J}{\sim}\xi$.
\SourceFigure{GELFAND-P000493}{dimension-margin}{Волнистое подчёркивание у пометы справа.}


Тогда $\underline{A\xi\overset{J}{\sim}\xi}$.

[DIAGRAM]
\SourceFigure{GELFAND-P000493}{spheres-label}{У слова «сферы» рамка с открытой левой стороной.}

\SourceFigure{GELFAND-P000493}{sphere-square}{Схема с двумя уровнями и вертикальными многоточиями.\par\smallskip
{\centering
[DIAGRAM]\par}}

\SourceFigure{GELFAND-P000493}{left-base-notation}{В левом нижнем узле написано $B(n+1)$.}

\Unclear{GELFAND-P000493}{right-sphere-arrow}{Правая верхняя стрелка сильно наведена. Остриё видно слева.}


Получ. послойная гомот. эквив.

\sout{Но: как их}
\SourceFigure{GELFAND-P000493}{question-cancellation}{Слова «Но: как их» зачёркнуты.}

Их нет, но как же их «найти»?
\SourceFigure{GELFAND-P000493}{j-factorization}{От $X$ к $BU(n)$ проведены две стрелки.\par\smallskip
{\centering
[DIAGRAM]\par}}

(\uline{классиф. пр-во рассл. со слоем}
$\overset{\text{\scriptsize (гомотоп.)}}{\text{\uline{сфера}}}$)
\SourceFigure{GELFAND-P000493}{homotopy-sphere-insert}{Перед «сфера» вписано «гомотоп.».}


$J$ — функтор.

Пусть $\exists BU(n)^{\wedge}$, на к-рое перен. операции Адамса,
т. е. $\exists$ опер. $(\psi^k)^{\wedge}$,
для к-рых коммут. диагр.
\SourceFigure{GELFAND-P000493}{completed-adams-diagram}{Правая нижняя стрелка изгибается вверх к $B_F$.\par\smallskip
{\centering
[DIAGRAM]\par}}

Быть может, может быть, что $(\psi^k)^{\wedge}$
\uline{сохр. размерность.}

(\uline{ориентации} — сохр. разм. отобр.)

Эти пр-ва постр. Сулливаном.


\Needspace{7\baselineskip}

\SourcePage{GELFAND-P000494}{unlabelled}
$BU(n)^{\wedge}$ будет пополн. грассман.

На
\Unclear{GELFAND-P000494}{initial-cancellation}{После «На» зачёркнута короткая запись. Чтение не установлено.}

\uuline{\mbox{$S^{2n-1}$} гр. \mbox{$U(n)$}}.
\Unclear{GELFAND-P000494}{sphere-exponent}{Знак в показателе густо наведён. Условно набрано $2n-1$. Возможно $2n+1$.}

Получ. расслоения (сфер.) на $BU(n)$.

\uline{Сулливан:} другое пр-во $\widehat S^{2n-1}$ и гр. $G_n$.
Рассм. рассл. и классиф. пр-во.
\SourceFigure{GELFAND-P000494}{top-completion-diagram}{Горизонтальная стрелка обведена снизу. Справа отдельно написано $\widehat S^{2n+1}$.\par\smallskip
{\centering
[DIAGRAM]\par}}

и будет сохр. послойный гомот. тип.

\Incomplete{GELFAND-P000494}{obliterated-line}{Строка густо перечёркнута. Связное чтение не установлено.}


Рассматриваются односвязные компл. с кон. порожд. гомот. гр.

\uline{Утв.}\quad
$\overset{\forall X}{\exists}$
\SourceFigure{GELFAND-P000494}{forall-insert}{Над $\exists$ вписано $\forall X$.}

\SourceFigure{GELFAND-P000494}{universal-triangle}{Верхняя стрелка обведена снизу. Две стрелки идут к нижнему узлу.\par\smallskip
{\centering
[DIAGRAM]\par}}

(функтор)

$\alpha(X)$ — преобр. тожд. функтора.

($\forall F$, \raisebox{.5ex}{\scriptsize (у к-ро)}
$\pi_i(F)$ конечны).
\SourceFigure{GELFAND-P000494}{finite-insert}{После $F$ вписано «у к-ро».}

\Unclear{GELFAND-P000494}{finite-cancellation}{Перед «конечны» затушёван фрагмент. Чтение не установлено.}


\uuline{\mbox{$\pi_i(\widehat X)$} проконечны}
($\underset{\longleftarrow}{\lim}$ кон. групп).
(в частн., компактны).

\sout{В частности}
\Unclear{GELFAND-P000494}{before-then-cancellation}{Условное чтение зачёркнутого «В частности».}

Тогда можно выч.:
\[
\pi_i(\widehat X)=\widehat{\pi_i(X)}
\]
(прокон. пополнение).
Окр. — подгруппы кон. индекса и их предел.

$Z^{\wedge}$ опис. через
$\overset{\text{\scriptsize (целые)}}{\text{\mbox{$p$}-адич. числа}}$.
\SourceFigure{GELFAND-P000494}{padic-insert}{Перед $p$ вписано «целые».}


Гом. группы $\widehat X$ ведут себя как целые $p$-ад. числа.

\sout{Действие гр. Галуа.}
\SourceFigure{GELFAND-P000494}{galois-cancellation}{Заголовок перечёркнут.}

\SourceFigure{GELFAND-P000494}{completion-square}{Нижняя стрелка направлена вправо, над ней знак равенства.\par\smallskip
{\centering
[DIAGRAM]\par}}

все гом. гр. конечны (это стаб. гомот. группы сфер)

\SourceFigure{GELFAND-P000494}{lower-completion-diagram}{Нижняя горизонтальная стрелка направлена вправо. У её левого конца виден дополнительный штрих. Стрелка от $\widehat S^{2n+1}$ проходит через соседнюю надпись.\par\smallskip
{\centering
[DIAGRAM]\par}}


\sout{Пусть}
\Unclear{GELFAND-P000494}{adams-cancellation}{Следующее короткое слово зачёркнуто. Чтение не установлено.}

\raisebox{.5ex}{\scriptsize (Постр. опер. $\psi_k$)}
прод. опер. Адамса и сохр. размерность.
\Unclear{GELFAND-P000494}{adams-operation-index}{Малый знак справа от $\psi$ условно набран как индекс $k$. Возможна запись $\widehat\psi^k$.}


Рассм. грассманы: $G_{N,n}$.
Это конечн. число афф. алг. карт, где ур-я над $Z$ (сх. над $Z$).
— алг. мн-зие над $Q$.
$\Rightarrow$ там действ. гр. Галуа $\operatorname{Gal}(\mathbb C:Q)$
— действует на $G_{N,n}$.
Оно не непр. в об. смысле, но непр. в топ. Зарисского.
Рассм. \uuline{Étale}-топ. \ldots
\[
\operatorname{Gal}(\mathbb C:Q)\longrightarrow
\operatorname{Gal}\bigl(
\underset{\substack{\vert\\\text{\scriptsize max аб. расш.}}}{Q_{ab}};Q\bigr)
=\widehat Z^{*}
\]
\SourceFigure{GELFAND-P000494}{abelian-label}{Под $Q_{ab}$ подписано «max аб. расш.».}

\[
H_i(V_c,A)\approx H_i(V_{et},A)
\]
\[
\widehat V_c\sim\underset{\longleftarrow}{\lim}\,N(V_{et}),
\qquad N\text{ — нерв}.
\]


\Needspace{7\baselineskip}

\SourcePage{GELFAND-P000495}{unlabelled}
\uuline{26.3.73}

\uwave{Формулир. рез.} \uwave{\mbox{\uuline{Атья, Патоди, Зингер.}}} \uwave{(И. М. Гельфанд).}

Пусть $A$ — самосопр. элл. опер.
\[
\zeta_A(s)=\operatorname{Tr}A^{-s}=\sum\lambda_i^{-s}
\]
при $s=0$ не имеет полюса (Сили).
\Unclear{GELFAND-P000495}{seeley-spelling}{Фамилия условно прочитана как «Сили». Возможно «Силли».}


$\zeta_A(0)$ — Атья и другие.

Рассм. Вместо $\zeta$-ф-и $L$-ряды.
\Incomplete{GELFAND-P000495}{eta-cancelled-phrase}{Зачёркнутая фраза справа, возможно, начинается с «Не». Продолжение не читается.}


\begin{flushright}\small \uline{Пусть \mbox{$A$} — не полож. опред.} (самосопр. эллипт.)\end{flushright}
\[
\eta(s)=\sum_{\lambda\ne0}(\operatorname{sign}\lambda_i)\,|\lambda_i|^{-s}
\]
Трудная теор.: \uline{\mbox{$s=0$} — не полюс} !!

Сущ. пример — опер. 1-го порядка.

Общая зад.: нет полож. опред. Рассм.
$\sum$ (мульт. характеров от $\lambda_i$)
($L$-ряды).
\uline{На компл. числах} \uline{быв. более сложные мульт. характеры.}

\uline{Пример.}
\Unclear{GELFAND-P000495}{example-cancellation}{Запись перед $A$ зачёркнута и частично налегает на $A$.}

\[
 A=i\frac{d}{dx}+t\qquad\text{на }[0,2\pi]\qquad\text{периодич.}
\]
\[
\underline{\eta(0)=1-2t}
\]
\SourceFigure{GELFAND-P000495}{eta-value-domain}{Формула подчёркнута скобкой, открытой сверху.}


Соб. знач. $n+t$, $n=-2,-1,0,1,\ldots$

$\eta(0)$ — «асимметрия» оператора:
\[
\sum_{\lambda_i>0}1-\sum_{\lambda_i<0}1.
\]

\uline{С. И. Гельфанд.} Излож. последней лекции Манина:

$p$-адич. интегрирование (Мазур, Суиннертон — Дайер)
\Unclear{GELFAND-P000495}{swinnerton-spelling}{Середина первого компонента фамилии написана слитно. Условное чтение «Суиннертон».}



\Needspace{7\baselineskip}

\SourcePage{GELFAND-P000496}{unlabelled}
$Z_p$ — целые $p$-адич. числа, на них задана ф-я. Надо проинтегрировать.

(знач. $p$-адические)

\uuline{лит.} Л. Г. Шнирельман, Дворк, Грауэрт, Краснер, Тэйт

$Z_p\supset A_i=p^iZ_p$ — сист. окр. $0$.
Мера зад. на мн-вах $a+A_i$

$\mu(a+A_i)$ — $p$-адич. число (не обяз. целое).
\[
a\in Z_p\quad\text{или}\quad \underline{a\in Z_p/p^i}
\]
Должна быть конечн. аддитивность; свод. к усл.:
\[
\mu(a+A_i)=\sum_{t=0}^{p-1}\mu(a+tp^i+A_{i+1})
\]

\uline{2 примера}: 1) \uline{инвар. мера}: не зав. от $a$, а зав. от $i$.
Тогда
\[
\mu(a+A_i)=p^{-i}\qquad\text{(единств. возможность).}
\]
Мера мал. мн-ва велика! — далеко от $0$
\[
\left(\left\|\frac ab\right\|=p^{-n},\quad
\text{где }\frac ab=p^n\frac{a_1}{b_1},\quad
a_1,b_1\text{ вз. пр. с }p\right).
\]

2) \uline{Более важный пример.} (меры Леопольда — Куботы).
\begin{flushright}
2 параметра $n$ и $a$
\end{flushright}
\[
(n,p)=1,\qquad 0\le a\le p^i-1
\]
\[
\mu_n(a+A_i)=\left[\frac{na}{p^i}\right]-\frac{n-1}{2}
\]
\[
\mu_2(a+A_i)=
\begin{cases}
+1/2,&\text{если }a>\dfrac{p^i}{2}\\
-1/2,&\text{если }a<\dfrac{p^i}{2}
\end{cases}
\]
Она аддит. !! Мера $\forall$ мн-ва ограничена

(она получ. преобр. Меллина из $\vert$ $\zeta$-ф-и Леопольда — Ивасавы)
\Unclear{GELFAND-P000496}{mellin-vertical-stroke}{Между «из» и $\zeta$ проведена вертикальная черта. Её назначение не установлено.}


\[
Q_p\ni f(x),\quad x\in Z_p\qquad
\int_{Z_p}f(x)\mu(x)\qquad\text{по Риману}
\]
$R_m$ — представители $Z_p/p^mZ_p$ (по 1 из $\forall$ класса)

Это разбиение. Сост.
\[
S(R_m)=\sum_{a\in R_m}f(a)\mu(a+A_m)
\]
они сходятся при усл. на меру и ф-ю.

\uline{Умеренно раст. мера} — если растёт медленнее инвариантной
\[
\lim_{m\to\infty}\max|\mu(a+A_m)|p^{-m}=0
\]
$f$ уд. усл. Липшица, если
\[
 |f(a)-f(b)|<Cp^{-m},\quad\text{если }|b-a|<p^{-m}
\]
(в частн., если зад. сход. степ. рядом)

Если мера умер. растёт, а ф-я удовл. усл. Липшица, то, $\exists$ предел, наз. интегралом.


\Needspace{7\baselineskip}
